\scanpage{30}
Or $\widetilde X^{\nu+1}$ est ind-représentable de la façon suivante. On regarde sur $X^{(\nu)}:=X^{\nu+1}$ les voisinages infinitésimaux $\Delta_{X/S}^{(\nu)}(i)$ de la diagonale. Pour chacun, on regarde (grâce à l'augmentation $\underline O_{\Delta^{(\nu)}(i)}\longrightarrow O_{\Delta^{(\nu)}(\unsure{0|\omega})}=\underline O_X$) le \unsure{fourbi|fonct.} augmenté à \uline{puissances divisées} qu'il définit, soit $\widetilde{\Delta_{X/S}^{(\nu)}(i)}$ [\deleted{\unsure{ainsi on}} c'est une notation \unsure{doublement} justifiée, i.e. cohérente avec \deleted{celle déjà} l'usage de $\sim$ déjà utilisé\ldots], et ses $j$-tronqués [\deleted{pour} \deleted{$j$} \deleted{le} au sens des puissances divisées] $\Delta_{X/S}^{(\nu)}(i,j)$, de sorte que du point de vue
\[
\underbrace{\widetilde{\Delta_{X/S}^{(\nu)}(i)}}_{\begin{gathered}
\text{faisceau}\\
\deleted{\text{De Rham}}\ \mathrm{DR}\text{-stratifiant}\\
\text{``\deleted{\unsure{ind}}représenté'' par }\Delta_{X/S}^{(\nu)}(i)\text{ augmenté}\\
\text{\unsure{vers}}\quad\Delta_{X/S}^{(\nu)}\deleted{\unsure{(i)}}(\unsure{0|\omega})
\end{gathered}}
=\varinjlim_j\Delta_{X/S}^{(\nu)}(i,j)
\]
et
\[
\widetilde X^{\unsure{2\nu+1|\nu+1}}=\varinjlim_{i,j}\widetilde{\Delta_{X/S}^{(\nu)}(i,j)}.
\]
On aura aussi
\[
H^q(\widetilde{\mathrm{DRst}}/\widetilde{(U,U')},O)\simeq H^q(U',\underline O_{U'})
\]
\margin{à vérifier}
\mspage{29}
